\documentclass[11pt,a4paper]{article}
\usepackage{amsmath,amssymb,booktabs,hyperref,geometry}
\geometry{margin=25mm}
\title{Flowed Clover Topological Charge: Normalization and Verification}
\author{PyQUDA Measurement}
\date{October 2026}
\begin{document}
\maketitle
\section{Scope and conventions}
This implementation measures only a hadron-independent gluon loop. It does
not extract a nucleon form factor or perform a continuum extrapolation.
Use directions $(0,1,2,3)=(x,y,z,t)$ and $\epsilon_{0123}=+1$.
Hermitian generators satisfy $\operatorname{Tr}(T^aT^b)=\delta^{ab}/2$.
For links $U_\mu=\exp(iag A_\mu)$ and the sum $C_{\mu\nu}$ of four oriented
plaquettes based at the same site, define
\begin{align}
 A_{\mu\nu}&=\frac{C_{\mu\nu}-C_{\mu\nu}^\dagger}{8}
 -\frac{\mathbf1}{3}\operatorname{Tr}\frac{C_{\mu\nu}-C_{\mu\nu}^\dagger}{8},\\
 H_{\mu\nu}&=-iA_{\mu\nu}=a^2gG_{\mu\nu}+O(a^4).
\end{align}
Here $A$ is anti-Hermitian, $H$ is Hermitian, and both are traceless.
In particular $H_{\nu\mu}=-H_{\mu\nu}$.

\section{The three-trace formula}
The 24 nonzero permutations in the epsilon contraction fall into three
pairings. Reversal within either antisymmetric pair cancels the corresponding
epsilon sign; exchanging the pairs leaves their trace unchanged by cyclicity.
Each pairing therefore contributes eight equal terms, giving
\begin{align}
 q_{\mathrm{lat}}
 &=\frac{\epsilon_{\mu\nu\rho\sigma}}{32\pi^2}
    \operatorname{ReTr}(H_{\mu\nu}H_{\rho\sigma})\\
 &=\frac1{4\pi^2}\operatorname{ReTr}
    (H_{01}H_{23}-H_{02}H_{13}+H_{03}H_{12}).
\end{align}
No commutativity assumption is needed. Since
$\operatorname{Tr}(HK)^*=\operatorname{Tr}(KH)=\operatorname{Tr}(HK)$,
the local density is real. It is invariant under local color conjugation and
changes sign under reversal of one spatial axis. For $H_{01}=H_{23}
=\operatorname{diag}(1,-1,0)/2$ and other independent planes zero,
$q_{\mathrm{lat}}=1/(8\pi^2)$, fixing a simple normalization test.

The links already contain the coupling:
\begin{equation}
 q_{\mathrm{lat}}=a^4q_{\mathrm{phys}}+\text{cutoff corrections},\qquad
 q_{\mathrm{phys}}=\frac{g^2}{32\pi^2}
 \epsilon_{\mu\nu\rho\sigma}\operatorname{Tr}(G_{\mu\nu}G_{\rho\sigma}).
\end{equation}
If $G$ includes $g$, omit the explicit $g^2$ in the last expression.
No additional flavor factor, fermion-loop sign or EMT factor of two belongs
in the implementation. QUDA's anti-Hermitian lower-plane convention maps to
this formula using $A=iH$; two factors give the required minus sign.

\section{Fourier transform and source alignment}
For $V_s=L_xL_yL_z$, origin $o=0$ and integer momentum $\mathbf n$, store
\begin{equation}
 L(\mathbf n,t_f,\tau)=\frac1{V_s}\sum_{\mathbf x}
 e^{2\pi i\sum_j n_j(x_j-o_j)/L_j}q_{\mathrm{lat}}(t_f,\mathbf x,\tau).
\end{equation}
Reality implies $L(-\mathbf n)=L(\mathbf n)^*$ and
$Q(t_f)=V_s\sum_\tau L(0,t_f,\tau)$.
For a source $s=(\mathbf s,t_0)$,
\begin{equation}
 L_s(\mathbf n,t_f,\tau)=e^{-2\pi i\sum_jn_j(s_j-o_j)/L_j}
 L(\mathbf n,t_f,(t_0+\tau)\bmod N_t).
\end{equation}
The gluonic time shift has no antiperiodic fermion sign. The disconnected
building block with a compatible polarized two-point contraction is
\begin{equation}
 C_{3,\Gamma}^{Q}=V_s\bigl(\langle C_{2,\Gamma}L_s\rangle
 -\langle C_{2,\Gamma}\rangle\langle L_s\rangle\bigr).
\end{equation}
Configuration/stream identity and source correlations must survive averaging.
This is not yet a complete kinematic ratio.

\section{Flowed observable and RGI limit}
The output interval and integration step are distinct. By default all intervals
use the same substep count. An optional first-interval override permits ten
substeps initially and three thereafter without changing the output times.
Both counts and step sizes are recorded in metadata; omitting the override
preserves the uniform integration schedule.
At fixed positive physical $t_f$, take the continuum limit before $t_f\to0$:
\begin{equation}
 F_Q^{\mathrm{RGI}}(Q^2)=\lim_{t_f\to0}\lim_{a\to0}
 F_Q^{\mathrm{GF}}(Q^2,a,t_f).
\end{equation}
The small-flow expansion identifies the RGI density, including the singlet
axial-divergence mixing structure~\cite{luscher,lw}. No separate axial
measurement or standalone $\overline{\mathrm{MS}}$ conversion is performed
here. Finite-flow results must not be labelled already RGI; running and mixing
can produce non-polynomial, logarithmic corrections in the extrapolation.
Global continuum charge conservation along flow does not imply independence
of local nonforward matrix elements. At finite spacing even the clover total
charge has cutoff effects. Require a physical window with
$a\ll\sqrt{8t_f}$ and a radius small relative to insertion/source separations.

In the proposed nucleon convention~\cite{fukushima},
\begin{equation}
 M_N^{-1}\langle N(p')|q_{\rm top}|N(p)\rangle
 =\bar u(p')i\gamma_5u(p)F_Q(Q^2),\quad d_Q=-F_Q(0).
\end{equation}
The forward pseudoscalar bilinear vanishes: extrapolate from finite transfer.
In general $Q^2=|\mathbf p'-\mathbf p|^2-(E'-E)^2$.
The Wick-rotation sign must be derived together with gamma matrices, epsilon
and the chosen polarized projector; the explicit Minkowski $i\gamma_5$ is not
an extra $i$ multiplying the real Euclidean density. A Breit-frame density is
not an ordinary relativistic probability density. Existing overlap/AWI
calculations~\cite{gong} guide signal construction, not direct flowed-clover
matching.

\section{Verification and cost}
CPU tests compare all 24 epsilon terms with the compact expression on
noncommuting matrices, test gauge conjugation, reflection, Hermiticity,
tracelessness, identity links, Fourier pairing, source alignment, scheduling,
root ownership and HDF5 metadata. GPU tests check actual link-level gauge
invariance, MPI decompositions, QUDA total charge and the pre-extraction EMT
baseline. Integration-step refinement compares equal physical output times;
finite-step outputs need not be identical. See the companion Markdown
validation summary for actual pass/pending status.
The installed QUDA exponential has a singularity at exactly zero force.
Identity links are therefore checked only at zero flow time; nontrivial
fields exercise flow. This restriction does not exclude zero Fourier momentum
or zero total charge, and does not constitute a fix of QUDA.

Six clover planes and three traces are needed per output time, with no quark
inversions. Only projected loops are transferred and stored. An output with
125 momenta, 21 flow times and 80 time slices occupies 3.20 MiB in complex128;
a full real $80^4$ density would occupy 312.5 MiB per flow time.

\subsection{Validated Aurora result}
Job 8948692 completed successfully on 20 October 2026 with 10 nodes and 100
ranks. Its real $80^4$ input was stream b, configuration 600, Coulomb gauge
fixed with no additional HYP for topology. Both Wilson and Symanzik flow used
27 momenta and output times $t_f/a^2=0,0.1,0.2,0.3,0.4$.
The pre-extraction EMT baseline was commit \texttt{5036ce0}; all compared EMT
datasets agreed exactly. Small-field MPI-decomposition and gauge-transformation
differences were below $1.92\times10^{-16}$ and $1.01\times10^{-15}$ respectively.
Total-charge differences relative to native QUDA were at most
$2.67\times10^{-13}$ (Wilson) and $1.75\times10^{-13}$ (Symanzik).
Step halving from $0.01$ to $0.005$ to $0.0025$ reduced consecutive loop
differences by factors 8.22 and 8.36 respectively. The identity-flow limitation
above remains in force; this is not a validation of that singular case.

Two clean repeats, excluding warmup and initial gauge input, took about
39.27 seconds per Wilson measurement and 41.02 seconds per Symanzik measurement.
Separately instrumented passes assigned 32.56/34.28 seconds to flow and
6.56/6.56 seconds to clover construction. Density contraction was about
0.063 seconds, projection 0.016 seconds and HDF5 output 0.011 seconds or less.
These are maximum accumulated component times over ranks, not strictly
additive wall times. The companion Markdown report retains detailed results.

\subsection{Production-schedule benchmark}
Job 8948902 also passed on 20 October 2026, using ten nodes and 100 ranks.
Five configurations (b.600--640 in steps of ten) were measured without HYP,
then with one HYP step. Symanzik output times were
$t_f/a^2=0,0.207936,\ldots,2.07936$, with ten integration substeps initially
and three per later interval, and 125 momenta including zero.
All ten files passed finite-value, Fourier and metadata checks. Maximum
total-charge differences from native QUDA were $1.60\times10^{-13}$ (noHYP)
and $8.66\times10^{-14}$ (HYP). This is not a step-refinement test at the
production step size.
Excluding the first configuration, production-like totals averaged
52.75 and 53.19 seconds respectively; independent native-reference work
added about 36.6 seconds and is excluded from those totals, as are
initialization and shutdown. HYP ran second, so input-cache effects may
affect the comparison. These timings do not establish large-job throughput.

\begin{thebibliography}{9}
\bibitem{fukushima} Fukushima and Uji, \emph{Polarized Nucleon as a
Topological Dipole}, \href{https://arxiv.org/abs/2607.03123}{arXiv:2607.03123}.
\bibitem{luscher} M. Luescher, \emph{On the chiral anomaly and the Yang--Mills
gradient flow}, \href{https://arxiv.org/abs/2109.07965}{arXiv:2109.07965}, Eq.~(21).
\bibitem{lw} M. Luescher and P. Weisz, \emph{Renormalization of the topological
charge density in QCD with dimensional regularization},
\href{https://arxiv.org/abs/2103.15440}{arXiv:2103.15440}, Sec.~6.
\bibitem{gong} M. Gong et al., \emph{Strange and Charm Quark Spins from Anomalous
Ward Identity}, \href{https://arxiv.org/abs/1511.03671}{arXiv:1511.03671}.
\end{thebibliography}
\end{document}
